\section{Typing and Metatheory}

\subsection{The Audit Judgment}

The core target judgment is:
\[
  \judgment.
\]
This judgment is best read as an audit contract for the generated sidecar.
Under capability context $\mathcal{K}$, quantum context $\Gamma$, and linear resource context $\Delta$, target program $t$ transforms the contexts to $\Gamma'$ and $\Delta'$ and reports effect $E$.
For the running example, the judgment says that the checker can follow each qubit as it leaves \code{SurfaceD5}, enters \code{Tetra15} for a hybrid region, returns to \code{SurfaceD5}, and consumes resource identifiers exactly once along the way.
The judgment is parameterized by a certificate policy $\pi$, so the same syntactic plan may be accepted in exploratory mode and rejected in strict mode if it uses \Assumed{} rules.

Figure~\ref{fig:typing} shows the rules that matter most for the story.
The direct rule captures why $\gate{H}$ on \code{SurfaceD5} is routine but direct $\gate{CNOT}$ is rejected: the current code index must provide a capability witness.
The switch rule captures why a dense region can temporarily use another code: the type of the qubit changes from $Q[C_1,d,\State]$ to $Q[C_2,d',\State]$.
The produce and consume rules capture why a magic state behaves linearly: once the identifier leaves $\Delta$, a second consume step cannot typecheck.

\begin{figure}[t]
\[
\frac{
  q:Q[C,d,\mu]\in\Gamma \quad \Cap(C,g)\in\mathcal{K} \quad \pi(\ell)=\mathsf{true}
}{
  \mathcal{K};\Gamma;\Delta \vdash \mathsf{direct}\ g\ q : \Gamma;\Delta \;!\; E_g
}
\]
\[
\frac{
  q:Q[C_1,d,\State]\in\Gamma \quad \Switch(C_1,C_2)\in\mathcal{K} \quad \pi(\ell)=\mathsf{true}
}{
  \mathcal{K};\Gamma;\Delta \vdash \mathsf{switch}\ q:C_1\to C_2 :
  \Gamma[q\mapsto Q[C_2,d',\State]];\Delta \;!\; E_{\mathsf{sw}}
}
\]
\[
\frac{
  r\notin\mathrm{dom}(\Delta) \quad \Factory(R)\in\mathcal{K}
}{
  \mathcal{K};\Gamma;\Delta \vdash \mathsf{produce}\ r:R :
  \Gamma;\Delta,r:R \;!\; E_{\mathsf{prod}}
}
\qquad
\frac{
  r:R\in\Delta
}{
  \mathcal{K};\Gamma;\Delta \vdash \mathsf{consume}\ r\ \mathsf{for}\ g\ q :
  \Gamma;\Delta-r \;!\; E_{\mathsf{use}}
}
\]
\caption{Representative typing rules for capability use, code switching, and linear probabilistic resources. The rules are chosen because each corresponds to a concrete checker behavior in the case study.}
\label{fig:typing}
\end{figure}

\subsection{What Is Mechanized Today}

The Lean development mechanizes the small core needed to make these checker invariants precise.
It defines codes, gates, modes, resource kinds, certificate levels, policies, primitive atoms, resource contexts, effects, and typed programs.
The most important theorem for the running example is the unsupported-direct-gate rejection theorem: if the capability predicate for $(C,g)$ is false, there is no well-typed direct step for $g$ on code $C$.
This theorem is the formal counterpart of the negative test that rejects direct $\gate{CNOT}$ on \code{SurfaceD5}.

The second important theorem family concerns resource linearity.
Lean proves that after a resource has been consumed, consuming the same identifier again is impossible in the resulting resource context.
The implementation mirrors this theorem with produced and consumed sets in the checker.
The point is not that the proof is mathematically deep; the point is that a common compiler error has a direct expression in the calculus, the checker, and the test suite.

The third theorem family concerns certificates and effects.
Lean proves that typed programs contain only certificate levels allowed by the selected policy, and that sequential and branch effect composition conservatively bounds component effects.
For the case study, this means a strict run rejects plans that depend on prototype bridge rules, and an exploratory run can still report the same plan with visible warnings.
It also means the total effect summary is not a detached table entry; it is the result of applying the same composition operators used in the formal core.

\subsection{Conditional Soundness Boundary}

The intended soundness statement has a local layer and a physical layer.
The local layer says that if $\mathcal{K};\Gamma;\Delta\vdash t:\Gamma';\Delta'!E$ under a strict policy, then every acquisition atom in $t$ is justified by an allowed rule, every linear resource is used consistently with $\Delta$, and the reported effect is a conservative composition of rule effects.
This layer is represented today by Lean, Z3, Python checker obligations, and geometry checks.

The physical layer is conditional on the imported theorem dependencies attached to \Certified{} rules.
For example, the Steane/tetrahedral code-switch certificate locally checks the GF(2) algebraic gauge-fixing core and records the cited flag-circuit single-fault-tolerance theorem as an imported dependency.
Similarly, the Reed--Muller 15-to-1 certificate locally checks triorthogonality and low-weight $Z$-error detection, while the constant-time surface-code implementation profile is imported from an architecture-level result.
The paper's claim is therefore not that \system proves every physical gadget from first principles.
The claim is that, once local and imported rule contracts are made explicit, the calculus composes them without hiding illegal gates, duplicated resources, backend mismatches, or assumption boundaries.
